\section{Introduction and statements of results}


%\Anote{Introduce partition results before density results}



%\Anote{I changed the organization of the paper a bit. Combine smaller sections into bigger sections and move proofs of partitions before the density results. We can reverse to the old way later if we want.}

%\textcolor{red}{Move the result about compact groups to the front. Focus on compact groups. Result on Z is treated as an application.}

This paper is the first one in a series in which we study Bohr sets in sumsets. The second paper in the series is titled ``Bohr sets in sumsets II: countable abelian groups'' \cite{Griesmer-Le-Le-countable}.
Let $G$ be an abelian topological group. For a finite set $\Lambda$ of characters (i.e. continuous homomorphisms from $G$ to $S^1 := \{ z \in \C: |z|=1\}$) and $\eta > 0$, the set 
\begin{equation*} \label{eq:bohr1}
B(\Lambda; \eta) := \{ x \in G : | \gamma(x)-1 | < \eta \textup{ for all } \gamma \in \Lambda\}
\end{equation*} 
is called a \textit{Bohr set} or a \emph{Bohr neighborhood of $0$}. We refer to $\eta$ as the \textit{radius} and $|\Lambda|$ as the \textit{rank} (or \textit{dimension}) of the Bohr set. The set $B(\Lambda; \eta)$ is also called a Bohr-$(|\Lambda|, \eta)$ set. 

If $A, B \subset G$, the sumset and difference set of $A$ and $B$ are $A\pm B: = \{a\pm b: a \in A, b \in B \}$. If $c \in \Z$, we define $cA:=\{ ca: a \in A\}$.
